% ============================================================
%  Paper 22: Neutrino Isocurvature and the S8 Power Damping
%  Target: JCAP Compilation (SDGFT Series)
% ============================================================
\documentclass[a4paper,11pt]{article}
\usepackage{jcappub}

\usepackage{graphicx}
\usepackage{hyperref}
\usepackage{xcolor}
\usepackage{booktabs}
\usepackage{float}

% ── SDGFT shared macros ──────────────────────────────────────
\usepackage{sdgft-macros}

% ── Document ─────────────────────────────────────────────────
\begin{document}

\title{Neutrino Isocurvature and the S8 Power Damping}

\author{David A. Besemer}
\affiliation{Independent Researcher}
\emailAdd{david.besemer@sdgft.org}

\abstract{
We solve the cosmic matter clustering tension (S8) by activating mixed neutrino-isocurvature initial conditions. We show that freezing the neutrino isocurvature fraction at the topological value beta_iso = 1/36 ~ 0.0278 damps the matter power spectrum on small scales, yielding S8 ~ 0.791, in agreement with DES surveys.
}

\arxivnumber{SDGFT-2026-22}

\maketitle

\section{Introduction}
The S8 tension measures the discrepancy in matter clustering between CMB and weak lensing surveys. We resolve this using neutrino isocurvature.

\section{Theoretical Formulation}
Here we formulate the core mathematical relations governing the system:
\begin{equation}
  \betaiso = \frac{1}{36} \approx 0.027778 \quad (\text{isocurvature power fraction})
\end{equation}
\begin{equation}
  S_8 = \sigma_8 \sqrt{\frac{\Omega_m}{0.3}} \approx 0.791 \quad (\text{matter clustering amplitude})
\end{equation}
\section{Isocurvature Power Damping}
The isocurvature fraction suppresses the growth of structure on small scales during the radiation-to-matter transition, lowering S8 without changing Planck's H0.

\section{Conclusion}
The isocurvature fraction beta_iso = 1/36 resolves the S8 tension, providing a unified explanation for late-time structure growth.



\begin{thebibliography}{99}
\bibitem{Goldstein_Paper01} D. A. Besemer, ``Axiomatic Foundations of SDGFT,'' SDGFT-2026-01 (in preparation).
\bibitem{Goldstein_Paper04} D. A. Besemer, ``Effective Spectral Dimension and the Banach Fixed-Point Theorem in SDGFT,'' SDGFT-2026-04.
\bibitem{Goldstein_Code} D. A. Besemer, \texttt{sdgft} Python package, \url{https://github.com/david-besemer/sdgft} (2026).
\end{thebibliography}

\end{document}
